% Start preamble
\documentclass[12pt,a4paper]{article}
\usepackage{geometry}
 \geometry{
 a4paper,
 total={170mm,257mm},
 left=20mm,
 top=20mm,
 }
\usepackage[utf8]{inputenc}
\usepackage[T1]{fontenc}
\usepackage[pdftex]{graphicx}
\graphicspath{{./}}
\usepackage{enumitem}
\usepackage{pdfpages}
\usepackage{hyperref}
\usepackage{tikz}
\usepackage{attachfile}
\usepackage{epstopdf}
\usepackage{array}
\usepackage{multirow}
\usepackage{multicol}
\usepackage{float}
%\usepackage[table]{xcolor,colorbl}
\setlength{\textwidth}{17cm}
\setlength{\oddsidemargin}{-0.5cm}
\setlength{\evensidemargin}{-0.5cm}
%\setlenght{\headsep}{0cm}
\setlength\parindent{0pt}
%\setlength{\extrarowheight}{3pt}
\usepackage{listings}
%\usepackage{xcolor}

\input{arduinoLanguage.tex}
%%%%%% Counting oppgaves %%%%%%
 \newcount\questnum \questnum=0
 \def\oppgave{
            \advance\questnum by 1
	    \ifthenelse{\questnum>0\AND \questnum<9}
	    {
                \vskip 1cm
		\textbf{Oppgave}\hskip 5pt\the\questnum \hfill \hfill(6p)
		\vskip 3pt
		\hrule
	\vskip 0.5cm}
	{
                \vskip 1cm
		\textbf{Oppgave}\hskip 5pt \the\questnum \hfill \hfill(12p)
		\vskip 3pt \hrule \vskip 0.5cm }

		}

\begin{document}
\title{Prøve Hydraulikk}
\author{Faglærer: Fred-Olav Mosdal 90507684\\}
\maketitle
|
\newpage
\oppgave{}%1
\vskip 2.5pt 
I hydrauliske anlegg har vi energitap, hvilken energiform går dette taper over til?
\vskip 5pt 
\begin{tikzpicture}
	\draw[step=0.5cm,gray!90,very thin]  grid (17,4) ;
\end{tikzpicture}

\oppgave{}%2
Hydraulikkpumper kan deles inn i kategorier på flere måter, hva er forskjellen på konstantpumper og variable pumper?
\vskip 2.5pt 
\begin{tikzpicture}
	\draw[step=0.5cm,gray!90,very thin]  grid (17,5) ;
\end{tikzpicture}
\oppgave{}%3
Hvordan er forholdet mellom $V_1$ og $V_2$?
\vskip 5pt 
$$\includegraphics[width=0.6\textwidth]{../output/noGPLimages/kontlov.png}$$
\vskip 5pt 
\begin{tikzpicture}
	\draw[step=0.5cm,gray!90,very thin]  grid (17,2);
\end{tikzpicture}
\vskip 2.5pt 
\newpage
\oppgave{}%4
\vskip 2.5pt 
Sett navn på hva de ulike styringsinnretninge betyr:
\begin{flushright}
$\includegraphics[width=0.7\textwidth]{../output/noGPLimages/betjening.png}$
\end{flushright}
\vskip 5pt 
\vskip 2.5pt 
\newpage
\oppgave{}%5
\vskip 2.5pt 
Hvorfor brukes det mest forstyrte trykkbegrensningsventiler?
\vskip 5pt 
\begin{tikzpicture}
	\draw[step=0.5cm,gray!90,very thin]  grid (17,10) ;
\end{tikzpicture}
\vskip 2.5pt 
\oppgave{}%6
\vskip 2.5pt 
Hvorfor bruker vi en trykkbegrensningsventil i hydrauliske anlegg?
\vskip 5pt 
\begin{tikzpicture}
	\draw[step=0.5cm,gray!90,very thin]  grid (17,10) ;
\end{tikzpicture}
\vskip 2.5pt 
\newpage
\oppgave{}%7
\vskip 2.5pt 
Tegn og forklar prinsippet for en tilbakeslagsventil?
\vskip 10cm 
\begin{tikzpicture}
	\draw[step=0.5cm,gray!90,very thin]  grid (17,13) ;
\end{tikzpicture}
\vskip 2.5pt 
\oppgave{}%8
\vskip 2.5pt 
En jekk er utstyrt med en pumpe der stempelet (S1) har et areal på 1 cm². Kraften vi bruker på pumpa, er 600 N. Arealet til stempelet på løftesylinderen (S2) er 20 cm²:
\begin{itemize}
	\item Hvor stor kraft får vi på løftesylinderen?
	\item Hvor mange kilo kan den løfte? 
\end{itemize}
\vskip 5pt 
\begin{tikzpicture}
	\draw[step=0.5cm,gray!90,very thin]  grid (17,20) ;
\end{tikzpicture}
\vskip 2.5pt 
\newpage
\oppgave{}%9
\vskip 1cm 
Beskrivelse av problemet En dobbeltvirkende sylinder brukes til å åpne og lukke en skottluke. Lukking skal skje uten rykk og skal foregå i en jevn, justerbar hastighet. Hastigheten justeres med en enveis strømningskontrollventil.
\vskip 1cm 
For å sikre at den tunge døren ikke trekker stempelstangen ut av sylinderen under lukkingen, må en trykkavlastningsventil brukes som en motstøtte.
\vskip 1cm 
$$\includegraphics[width=1\textwidth]{../output/noGPLimages/a2INTHydraulikk02.png}$$
\vskip 1cm 
Du skal simulere kretsen i fluidsim
\vskip 1cm 
$$\includegraphics[height=1\textheight]{../output/noGPLimages/a2INTHydraulikk01.png}$$
\vskip 5pt 
\vskip 2.5pt 
\end{document}
